\textbf{The analytic Compton total cross section at next-to-leading
order is already known from the work of Roman~N.~Lee, Matthew~D.~Schwartz and
Xiaoyuan Zhang~\cite{qed-processes-lsz}. Our purpose is an independent
recalculation and verification of their result.}
We derive the virtual and real amplitudes from the Feynman rules and
integrate them with exact electron-mass dependence. The resulting
expression in logarithms and polylogarithms contains no unevaluated
loop-parameter or phase-space integral. This chapter states the
observable, conventions, and complete result before the detailed
derivation. The published coefficients are used only for the final
comparison, never as boundary data or inputs to the integration.

\section{Observable and conventions}
We consider unpolarized scattering on a free electron,
\begin{equation}
 e^-(p)+\gamma(k)\longrightarrow e^-(p')+\gamma(k'),
 \qquad
 e^-(p)+\gamma(k)\longrightarrow e^-(p')+\gamma(k_1)+\gamma(k_2).
\end{equation}
The electron pole mass is $m$, its charge is $-e_{\mathrm T}$, and
$\alpha_{\mathrm T}=e_{\mathrm T}^2/(4\pi)$ is defined in the Thomson
limit. We average over the two incoming electron spins and two photon
polarizations and sum over all final polarizations. All allowed photon
energies and directions are integrated. The distinct pair-production
channel $e^-\gamma\to e^-e^-e^+$ is excluded. This is the same
observable as the photonic total of Reference~\cite{qed-processes-lsz}.
The massive polarized calculation of Denner and Dittmaier also treats
pair production separately~\cite{qed-processes-denner}.

We use natural units $\hbar=c=1$ and the metric
$\eta_{\mu\nu}=\operatorname{diag}(1,-1,-1,-1)$.
Define the physical invariants and the
rest-frame photon energy $\omega$ by
\begin{align}
 s_{\mathrm{phys}}&=(p+k)^2=m^2+2m\omega,\qquad
 t=(p-p')^2,\qquad u=(p-k')^2,\\
 X&=\frac{s_{\mathrm{phys}}}{m^2}-1=\frac{2\omega}{m}>0,\qquad
 Y=\sqrt{\frac{X}{X+4}},\qquad
 q=\frac{\sqrt{X+4}-\sqrt X}{\sqrt{X+4}+\sqrt X}.
 \label{eq:comp-overview-variables}
\end{align}
The variable $q$ is unrelated to electric charge. All displayed square
roots are positive. For the two-body process,
$s_{\mathrm{phys}}+t+u=2m^2$.

Ultraviolet integrals are regulated in $d=4-2\epsilon_{\rm UV}$
dimensions with scale $\mu$. A photon mass $0<\lambda\ll m$ regulates
virtual soft momenta and the matching real soft integral. In Feynman
gauge the propagators and electron vertex are
\begin{equation}
 S_F(P)=\frac{\ii(\slashed P+m)}{P^2-m^2+\ii0},\qquad
 D_{\rho\sigma}(\ell)=\frac{-\ii\eta_{\rho\sigma}}
 {\ell^2-\lambda^2+\ii0},\qquad
 +\ii\mu^{\epsilon_{\rm UV}}e_{\mathrm T}\gamma^\rho.
\end{equation}
Here $\slashed P=\gamma^\mu P_\mu$, with
$\{\gamma^\mu,\gamma^\nu\}=2\eta^{\mu\nu}$.
The photon symbol omits $F$ but uses the same Feynman boundary
conditions. External hard photons are massless. Terms vanishing as
$\lambda/m\to0$ are omitted only after the singular terms have been
identified. The real and virtual contributions are combined before
removing the infrared regulator.

\paragraph{Dimensional and integration variables.} Chapters 1 and 2 use
dimensional momenta and photon mass. In the integration chapters,
$s=s_{\mathrm{phys}}/m^2$, $t=t_{\mathrm{phys}}/m^2$,
$u=u_{\mathrm{phys}}/m^2$, and $\lambda=\lambda_{\mathrm{phys}}/m$.
The external energy variables satisfy $H=X=s-1$. In formulas for
elastic angular integration, $c$ denotes the center-of-mass photon
scattering cosine. The laboratory angle is denoted by $\theta$. Lowercase integration
coordinates are defined locally and are not identified with $X$.
The final evaluator uses $C_1(X)$, while the sector assembly uses
$C_{\mathrm{total}}(s)=C_1(s-1)$, in both cases multiplying
$\alpha_{\mathrm T}^3/m^2$. The labels $UV$, $UW$, and $VW$ in the
real-emission calculation name pairs of angular denominators.
They do not denote ultraviolet counterterms.

\section{On-shell renormalization constants}
The renormalized electron propagator has its pole at $m$ with unit
residue. The transverse photon residue is normalized at zero momentum
squared, and the charge is fixed by the zero-momentum-transfer electron
vertex. These are the QED on-shell conditions of the general formalism
reviewed in Section 3 of Denner~\cite{compton-overview-denner-renormalization}.
For one electron species, use the bare-to-renormalized definitions
\begin{align}
 \psi_0&=Z_2^{1/2}\psi,\qquad A_{0\mu}=Z_3^{1/2}A_\mu,\qquad
 m_0=Z_m m,\\
 e_0&=\mu^{\epsilon_{\rm UV}}Z_e e_{\mathrm T},\qquad
 Z_1=Z_eZ_2Z_3^{1/2},\qquad Z_j=1+\delta_j+\Order(\alpha_{\mathrm T}^2).
\end{align}
Here $Z_1$ multiplies the electron--photon interaction vertex.
Thus $Z_m$ relates the bare mass to the pole mass, while the mass coefficient in the
Lagrangian is $Z_2Z_m$. Define
\begin{equation}
 U_m=\frac1{\epsilon_{\rm UV}}-\gamma_E+\ln4\pi
                   +\ln\frac{\mu^2}{m^2},\qquad
 C=\frac{\alpha_{\mathrm T}}{4\pi}.
\end{equation}
Here $\gamma_E$ is Euler's constant. In the photon-mass convention the
one-loop constants are
\begin{align}
 \delta_m&=-C(3U_m+4)+\alpha_{\mathrm T}o(1),\\
 \delta_2&=-C\left[U_m+4+2\ln\frac{\lambda^2}{m^2}\right]
                           +\alpha_{\mathrm T}o(1),\\
 \delta_1&=\delta_2,\qquad
 \delta_3=-\frac{\alpha_{\mathrm T}}{3\pi}U_m,\qquad
 \delta_e=\frac{\alpha_{\mathrm T}}{6\pi}U_m.
 \label{eq:comp-overview-z}
\end{align}
The $o(1)$ terms vanish as $\lambda/m\to0$ at fixed $\mu/m$.
The Ward identity gives $Z_1=Z_2$ and hence $Z_e=Z_3^{-1/2}$.
The field constants above refer to Feynman gauge. They are not
separately observable quantities.
These constants agree with the one-electron QED specialization of
Denner's on-shell formulas, Equations (3.19), (3.20), and (3.32) of
Reference~\cite{compton-overview-denner-renormalization}.
His photon field counterterm $\delta Z_{AA}$ is our $\delta_3$,
and in QED his charge relation reduces to $\delta_e=-\delta_3/2$.

The signs can be checked directly against the subtraction constants
used later. With $\Sigma(P)=a(P^2)\slashed P+mb(P^2)$, let
\begin{equation}
 c_m=m[a(m^2)+b(m^2)],\qquad
 c_Z=a(m^2)+2m^2[a'(m^2)+b'(m^2)].
\end{equation}
Primes differentiate the squared argument. The counterterm insertion
is $\ii[\delta_2(\slashed P-m)-m\delta_m]$, so
\begin{equation}
 m\delta_m=-c_m,\qquad \delta_2=c_Z,\qquad
 \Sigma_R=\Sigma-c_m-(\slashed P-m)c_Z.
 \label{eq:comp-overview-subtractions}
\end{equation}
In this convention the subtraction constants are
$c_m/m=C(3U_m+4)+\alpha_{\mathrm T}o(1)$ and
$c_Z=-C[U_m+4+2\ln(\lambda^2/m^2)]+\alpha_{\mathrm T}o(1)$.
The vertex counterterm is $c_Z\gamma^\rho$ after factoring out the
charge. No further external-electron residue is applied to amplitudes
already normalized this way. For the two external photons the charge
and photon factors obey $Z_e^2Z_3=1$. There is no internal Born photon
on which to insert vacuum polarization in Compton scattering.

\section{The complete total cross section}
Let $\mathcal M_0$, $\mathcal M_1$, and $\mathcal M_{\gamma\gamma}$
denote the Born, renormalized one-loop, and double-Compton amplitudes.
Their contributions have the normalization
\begin{equation}
 \begin{aligned}
 \sigma_0+\sigma_1=\frac1{4p\cdot k}\lim_{\lambda\to0}\biggl\{&
 \int\dd\Phi_2\left[\overline{|\mathcal M_0|^2}
       +2\operatorname{Re}\overline{\mathcal M_0^*\mathcal M_1}\right]\\
 &+\frac1{2!}\int\dd\Phi_3\overline{|\mathcal M_{\gamma\gamma}|^2}
 \biggr\}.
 \end{aligned}
\end{equation}
The bar denotes the spin and polarization average and sum specified
above, $\dd\Phi_n$ is Lorentz-invariant $n$-body phase space, and the
real soft region uses the same photon-mass convention as the virtual
term. The factor $1/2!$ accounts for the identical final photons.
For the coefficient $A(s)$ of the virtual infrared logarithm, with
$s=1+X=s_{\mathrm{phys}}/m^2$,
\begin{equation}
 \frac{m^2\sigma_{1,V}}{\alpha_{\mathrm T}^3}
 =A(1+X)\ln\frac\lambda m+V_{\rm fin}(X),\qquad
 \frac{m^2\sigma_{1,R}}{\alpha_{\mathrm T}^3}
 =-A(1+X)\ln\frac\lambda m+R_{\rm fin}(X).
\end{equation}
Here $V_{\rm fin}$ and $R_{\rm fin}$ are finite coefficients in the
same convention, and terms vanishing with $\lambda$ are understood.
The finite terms refer to a common photon-mass regulator. Their sum has the form
\begin{equation}
 \boxed{\sigma_{\mathrm{inc}}^{\mathrm{NLO}}(X)
 =\frac{\pi\alpha_{\mathrm T}^2}{m^2}B(X)
  +\frac{\alpha_{\mathrm T}^3}{m^2}C_1(X)
  +\frac1{m^2}O(\alpha_{\mathrm T}^4).}
 \label{eq:comp-overview-total}
\end{equation}
The subscript $\rm inc$ denotes the sum over the one-photon and
two-photon final states specified above. The remainder is at fixed
$X>0$. Put $r_\omega=X/2$. The Born coefficient is
\begin{equation}
 \begin{aligned}
 B(X)=2\bigg[&\frac{1+r_\omega}{r_\omega^2}
 \left(\frac{2(1+r_\omega)}{1+2r_\omega}
       -\frac{\ln(1+2r_\omega)}{r_\omega}\right)\\
 &+\frac{\ln(1+2r_\omega)}{2r_\omega}
       -\frac{1+3r_\omega}{(1+2r_\omega)^2}\bigg].
 \end{aligned}
 \label{eq:comp-overview-born}
\end{equation}
This is the integrated Klein--Nishina result, also written in
Equation (5) of Reference~\cite{qed-processes-lsz}.

To state the entire NLO coefficient, let $w$ denote an ordered list
of kernel labels, called a word, and let $\varnothing$ denote the
empty list. Define the generalized polylogarithms (GPLs) recursively by
\begin{align}
 G(a,w;X)&=\int_0^X\dd t\,f_a(t)G(w;t),\qquad G(\varnothing;X)=1,\\
 f_a(t)&=\frac1{t-a}\quad(a=0,1,-1,-2,-4),\qquad
 f_r(t)=\frac1{\sqrt{t(t+4)}},\\
 G(\underbrace{0,\ldots,0}_{n};X)&=\frac{\ln^nX}{n!}.
\end{align}
The last line fixes the logarithmically singular all-zero case.
The length of a word is its weight. Products of these integrals obey
shuffle identities: interleave the two lists in all ways preserving
the order within each list, counting multiplicities. This follows by
partitioning the product of their ordered integration domains.
The letter $r$ labels a square-root kernel, not a number to substitute
for $a$. The integrals follow the positive real axis. In the words
containing $1$, the apparent singularity is removable because the
inner factor is $\ln t/(t-1)$, with limit one at $t=1$.
All other nonzero letters lie outside the integration path.
Classical polylogarithms are defined by
$\operatorname{Li}_n(z)=\sum_{k=1}^{\infty}z^k/k^n$ for $|z|<1$
and continuation elsewhere. Their real boundary values below agree
with these real iterated integrals.

\begingroup
\def\comptonproductlabel{eq:comp-overview-product}
Products of elementary logarithms can make the collected result shorter without introducing additional polylogarithms. Define $L=\ln(1+X)$ and $M=\ln X$. The two GPLs with the same coefficient combine into
\begin{equation}
\mathcal P(X)=G(0,-1,0;X)+G(0,1,0;X)=\tfrac12\left[\operatorname{Re}\operatorname{Li}_3(X^2)-M\operatorname{Re}\operatorname{Li}_2(X^2)\right].
\end{equation}
For the derivation, first take $0<X<1$. The shuffle relation $G(0)G(0,a)=2G(0,0,a)+G(0,a,0)$, with $G(0,a;X)=-\operatorname{Li}_2(X/a)$ and $G(0,0,a;X)=-\operatorname{Li}_3(X/a)$, gives
\begin{equation}
 G(0,a,0;X)=2\operatorname{Li}_3(X/a)-M\operatorname{Li}_2(X/a),\qquad a=\pm1.
\end{equation}
Summing the two signs and using $\operatorname{Li}_n(z)+\operatorname{Li}_n(-z)=2^{1-n}\operatorname{Li}_n(z^2)$ gives the stated result. Both shuffle multiplicities and this duplication factor are needed for its factor $1/2$. For $X>1$, continue all functions on the same edge of the principal cut. On either edge, the imaginary parts of $\operatorname{Li}_2(X)$ and $\operatorname{Li}_3(X)$ are respectively $\pm\pi\ln X$ and $\pm\pi\ln^2X/2$, with the same sign. These follow by integrating their derivatives from $X=1$, and they cancel in $2\operatorname{Li}_3(X)-M\operatorname{Li}_2(X)$. The analogous cancellation holds for the paired expression at argument $X^2$. Thus both edges give the same real function and the continuation adds no branch constant. The following equivalent form uses only real polylogarithms with arguments in $(0,1]$:
\begin{equation}
\mathcal P(X)=\begin{cases}\tfrac12[\operatorname{Li}_3(X^2)-M\operatorname{Li}_2(X^2)],&0<X\leq1,\\ \tfrac12[\operatorname{Li}_3(X^{-2})+M\operatorname{Li}_2(X^{-2})]+\tfrac13M^3+\tfrac16\pi^2M,&X\geq1.\end{cases}
\end{equation}
To derive the second line, put $t=\ln X>0$, $D(t)=\operatorname{Li}_2(e^{-2t})$ and $T(t)=\operatorname{Li}_3(e^{-2t})$. Their derivatives are $D^{\prime}(t)=2\ln(1-e^{-2t})$ and $T^{\prime}(t)=-2D(t)$. The identity $\ln(e^{2t}-1)=2t+\ln(1-e^{-2t})$ then gives
\begin{align}
\operatorname{Re}\operatorname{Li}_2(e^{2t})&=\frac{\pi^2}{3}-2t^2-D(t),\\
\operatorname{Re}\operatorname{Li}_3(e^{2t})&=T(t)-\frac43t^3+\frac{2\pi^2}{3}t.
\end{align}
Indeed, the derivatives of the two sides are respectively $-2\ln(e^{2t}-1)$ and $2\operatorname{Re}\operatorname{Li}_2(e^{2t})$. Write $\zeta(n)=\sum_{k=1}^{\infty}k^{-n}$ for the Riemann zeta function at integers $n>1$. Their values at $t=0$ agree because $\operatorname{Li}_2(1)=\pi^2/6$ and $\operatorname{Li}_3(1)=\zeta(3)$. These derivative and boundary conditions fix both identities and give $\mathcal P(1)=\zeta(3)/2$.
The sole apparent rational pole at $X=1$ can also be removed before evaluation. Define the continuous elementary function
\begin{equation}
R_{\log}(X)=\begin{cases}\ln X/(X-1),&X\ne1,\\1,&X=1,\end{cases}\qquad R_{\log}(1+h)=1-\frac h2+\frac{h^2}{3}+O(h^3).
\end{equation}
The series follows by dividing the Taylor series of $\ln(1+h)$ by $h$, which converges for $|h|<1$. Extracting the common power $X^{-3}$ gives the exact total coefficient as
\begin{equation}
C_1(X)=\frac{1}{X^3}\left[R_{\log}(X)+\mathcal E(X,L,M)+\sum_{j=1}^{13}b_j(X,Y,L,M)\mathcal F_j(X)\right].
\label{\comptonproductlabel}\end{equation}
The elementary part is
\begin{equation}
\mathcal E=e_{0}L^{3}+e_{1}L^{2}M+e_{2}L^{2}+e_{3}LM+e_{4}L+e_{5}M+e_{6}.
\end{equation}
\begin{equation}
e_{0}=- \frac{X^{2} - 16 X - 23}{3}.
\end{equation}
\begin{equation}
e_{1}=4 X^{2} - 2 X - 12.
\end{equation}
\begin{equation}
e_{2}=- 2 X^{3} - \frac{11 X^{2}}{2} + \frac{133 X}{2} + 36 + \frac{96}{X + 4} - \frac{6}{X + 1} + \frac{1}{\left(X + 1\right)^{2}}.
\end{equation}
\begin{equation}
e_{3}=X^{2} + 5 X - 10 + \frac{12}{X + 1} - \frac{4}{\left(X + 1\right)^{2}}.
\end{equation}
\begin{equation}
e_{4}=\frac{37 X^{2}}{4} - 32 X - 30 - \frac{15}{X + 1} - \frac{1}{\left(X + 1\right)^{2}}.
\end{equation}
\begin{equation}
e_{5}=- \frac{14 X^{5} + 31 X^{4} + 14 X^{3} - 6 X^{2} - 8 X - 2}{2 \left(X + 1\right)^{3}}.
\end{equation}
\begin{equation}
e_{6}=- \frac{X \left(273 X^{3} - 982 X^{2} - 2960 X - 1744\right)}{24 \left(X + 1\right)^{2}}.
\end{equation}
The thirteen remaining functions and their coefficients are
\begin{equation}
\mathcal F_{1}=G(0,-1;X).
\end{equation}
\begin{equation}
b_{1}=X^{3} + X^{2} - 43 X + 6 - \frac{96}{X + 4} - \frac{6}{X + 1} + \frac{2}{\left(X + 1\right)^{2}}- 2 L X \left(2 X + 5\right).
\end{equation}
\begin{equation}
\mathcal F_{2}=G(-2,-1;X).
\end{equation}
\begin{equation}
b_{2}=- 4 X^{2} - 80 X - 68 - \frac{128}{X + 4} + \frac{8}{X + 1} - \frac{4}{\left(X + 1\right)^{2}}.
\end{equation}
\begin{equation}
\mathcal F_{3}=G(1,0;X),\qquad b_{3}=- X^{2} + X - 2.
\end{equation}
\begin{equation}
\mathcal F_{4}=G(r,-1;X).
\end{equation}
\begin{equation}
b_{4}=Y \left(3 X^{3} + 18 X^{2} + 44 X - 8 - \frac{64}{X}\right)- 3 M Y \left(X - 2\right) \left(X + 4\right).
\end{equation}
\begin{equation}
\mathcal F_{5}=G(0,-1,-1;X),\qquad b_{5}=- 5 X^{2} - 8 X + 32.
\end{equation}
\begin{equation}
\mathcal F_{6}=G(-1,-2,-1;X),\qquad b_{6}=8 \left(X^{2} - 4 X - 6\right).
\end{equation}
\begin{equation}
\mathcal F_{7}=G(-4,r,-1;X),\qquad b_{7}=Y \left(- 6 X^{2} + 72 + \frac{48}{X}\right).
\end{equation}
\begin{equation}
\mathcal F_{8}=\mathcal P(X),\qquad b_{8}=2 X^{2} + 4 X - 12.
\end{equation}
\begin{equation}
\mathcal F_{9}=G(r,0,-1;X),\qquad b_{9}=Y \left(- 5 X^{2} + 14 X + 88 + \frac{48}{X + 4} + \frac{48}{X}\right).
\end{equation}
\begin{equation}
\mathcal F_{10}=G(r,-1,0;X),\qquad b_{10}=- 3 Y \left(X - 2\right) \left(X + 4\right).
\end{equation}
\begin{equation}
\mathcal F_{11}=G(r,-1,-1;X),\qquad b_{11}=Y \left(15 X^{2} - 18 X - 216 - \frac{96}{X + 4} - \frac{96}{X}\right).
\end{equation}
\begin{equation}
\mathcal F_{12}=G(r,-2,-1;X),\qquad b_{12}=Y \left(- 16 X^{2} + 32 X + 256 + \frac{64}{X + 4} + \frac{192}{X}\right).
\end{equation}
\begin{equation}
\mathcal F_{13}=G(r,r,-1;X),\qquad b_{13}=9 X^{2} - 24.
\end{equation}

\endgroup
At $X=1$ the elementary quotient $\ln X/(X-1)$ is assigned its
continuous value one. These definitions specify the full physical
function for every $X>0$.
Equation~\eqref{eq:comp-overview-product} follows from our integrated
sectors by shuffle identities and rational coefficient collection.
After conversion to a common basis, it agrees coefficient by coefficient
with Equation (6) of Lee, Schwartz and Zhang~\cite{qed-processes-lsz}.
This equality is a result of the recalculation, not an input to it.

\section{Asymptotic limits}
The following limits are taken in the complete infrared-finite result.
For $X\to0^+$,
\begin{equation}
 \sigma_0=\frac{\pi\alpha_{\mathrm T}^2}{m^2}
 \left[\frac83-\frac83X+\frac{52}{15}X^2-\frac{133}{30}X^3
 +\Order(X^4)\right],
\end{equation}
\begin{equation}
 \begin{aligned}
 \sigma_1=\frac{\alpha_{\mathrm T}^3}{m^2}\biggl\{&
 X^2\left[-\frac{16}{9}\ln X+\frac7{15}\right]
 +X^3\left[\frac43\ln X-\frac{113}{90}\right]\\
 &+X^4\left[-\frac{25}{9}\ln X+\frac{203}{125}\right]
 +\Order\!\left(X^5|\ln X|\right)\biggr\}.
 \end{aligned}
 \label{eq:comp-overview-low}
\end{equation}
The correction vanishes, leaving the Thomson cross section
$8\pi\alpha_{\mathrm T}^2/(3m^2)$. This is consistent with Thirring's
low-energy theorem, whose all-orders formulation and generalizations
are discussed by Dittmaier~\cite{compton-overview-thirring}.
The $X^2\ln X$ and $X^2$ coefficients agree with Equation (10) of
Lee, Schwartz and Zhang~\cite{qed-processes-lsz}. The two further orders
shown here follow from the full expression stated above.

For $X\to\infty$, write $\ell=\ln X$. Then
\begin{equation}
 \sigma_0=\frac{\pi\alpha_{\mathrm T}^2}{m^2}
 \left[\frac{2\ell+1}{X}+\frac{18-8\ell}{X^2}
 +\Order\!\left(\frac\ell{X^3}\right)\right],
\end{equation}
\begin{equation}
 \sigma_1=\frac{\alpha_{\mathrm T}^3}{m^2}\biggl\{
 \frac1X\left[\frac{\ell^3}{3}-\frac{\ell^2}{2}
 +\frac{17\ell}{4}-\frac{75}{8}-\frac{\pi^2}{2}+4\zeta(3)\right]
 +\Order\!\left(\frac{\ell^3}{X^2}\right)\biggr\}.
 \label{eq:comp-overview-high}
\end{equation}
Here $\zeta(3)=\sum_{n=1}^{\infty}n^{-3}$.
All logarithmic coefficients and the finite constant agree with
Equation (11) of Reference~\cite{qed-processes-lsz}.
Replacing $X$ by $s_{\mathrm{phys}}/m^2=1+X$ in its leading high-energy
term changes only the stated remainder. The finite constant is fixed by the endpoint values of our integrated GPLs.
In particular,
\begin{equation}
 \frac{\sigma_1}{\sigma_0}
 =\frac{\alpha_{\mathrm T}}{6\pi}\ell^2
 +\Order(\alpha_{\mathrm T}\ell).
\end{equation}
Infrared cancellation therefore leaves an energy-dependent radiative
correction. These are fixed-order results, with no resummation of the
large energy logarithms.
